\workbookchapter{模型服务架构}

\begin{exercises}
\begin{exercise} 同一权重修改模板后是否仍是同一服务修订？分别从模型计算与接口语义解释。

\begin{solution}
权重相同只说明参数未改；模板改变特殊符号、角色边界和输入词元，从而改变$p(y\mid x)$的实际条件。因此它应形成新的服务制品修订，并冻结模板、分词器和生成配置。可保留共同权重哈希用于去重，但接口兼容与质量回归需重新判断。
\end{solution}
\end{exercise}

\begin{exercise} 为取消与完成同时发生的请求设计状态更新规则，并说明资源释放为何还需要独立幂等保护。

\begin{solution}
只允许非终态经CAS写入一个终态，取消和完成中先提交者获胜，另一方读回确定状态。资源可能由异步清理者释放，故还须以请求/分配代次的唯一释放记录保护引用递减；终态CAS只约束状态字段，不能自动覆盖另外的内存操作。迟到完成不得覆盖已取消状态，但远端业务结果仍需审计。
\end{solution}
\end{exercise}

\begin{exercise} 一条流已输出十个事件后连接中断，列出支持恢复所需的状态；说明事件重放与重新生成的区别。

\begin{solution}
至少保留请求身份、制品与采样配置、已提交事件日志和单调序号、终态、恢复保留期限与身份授权。若客户端已确认事件10，重连从11继续；重复事件按序号去重。事件重放读取已形成的结果，重新生成是新的随机计算，可能改变后文且重复成本，不能不加标记地冒充断点续传。
\end{solution}
\end{exercise}

\begin{exercise} 请求总预算2秒，排队已耗费1.2秒、网络与预处理耗费0.3秒，计算剩余预算并讨论是否应再发起一次1秒重试。

\begin{solution}
剩余预算 $2-1.2-0.3=0.5$ 秒。再发起预计1秒的尝试不能满足原期限，应拒绝重试、缩小为可兑现的子操作，或返回期限不足。所有下游必须接收原绝对截止时间或等价剩余预算；每重试重置2秒会破坏端到端约束。
\end{solution}
\end{exercise}

\begin{exercise} 区分启动、存活与就绪信号，说明过载时错误使用存活重启会发生什么。

\begin{solution}
启动信号允许加载和预热尚未完成；存活判断进程是否需要重建；就绪判断能否接收新请求。过载通常应降低准入或就绪，不能把暂时排队当死锁反复重启。重启会损失KV并引入冷启动，减少有效容量，形成更严重的过载反馈。
\end{solution}
\end{exercise}

\begin{exercise}\optional 对2 GiB缓存、40 GiB/s有效带宽，计算传输下界；每秒30请求时总带宽需求是多少？

\begin{solution}
纯传输下界 $\frac{2}{40}=0.05$ 秒，即50毫秒；30请求/秒要求60 GiB/s，超过单条40 GiB/s链路。该下界不含握手、布局、排队与所有权提交。需要增加有效带宽、减少迁移字节或限制接纳，不能因单请求50毫秒就认定系统可承载目标吞吐。
\end{solution}
\end{exercise}

\begin{exercise}\optional 在 P/D 迁移中，接收端已获得数据而所有权提交结果未知，应怎样避免双重输出？

\begin{solution}
接收端先处于已准备但无输出权状态。使用唯一请求代次和原子所有权记录确认提交结果；未知时查询权威记录并等待，不能双方凭超时自行继续。只有持当前代次租约/隔离令牌的一端发布结果，旧拥有者被执行或输出网关拒绝。数据复制成功不等于所有权迁移成功。
\end{solution}
\end{exercise}

\begin{exercise} 将容量算例启动时间改为120秒，计算新增积压，并提出两种不依赖即时扩容的处理方式。

\begin{solution}
现有容量1800、需求2400请求/分钟，恒定缺口600/分钟，120秒产生1200请求积压。两种可行方法是预热容量池/预测提前扩容，以及受期限约束的准入和背压；也可降级到经验证的低成本路径。若采用拒绝，积压公式应减去拒绝量，不能同时把拒绝请求算作仍排队。
\end{solution}
\end{exercise}

\begin{exercise} 比较影子、金丝雀与蓝绿方式的资源和证据边界，解释影子工具副作用为何需要隔离。

\begin{solution}
影子复制请求而不影响正式响应，消耗额外资源，验证输出差异但不能完全模拟真实用户反馈；金丝雀让部分真实请求使用候选，需随机或明确定义分流及回滚；蓝绿保留完整两套环境，切换快速但容量高。影子调用写工具会产生真实副作用，应只读、禁写或用明确隔离的业务验收环境，不能把“影子”当作副作用许可。
\end{solution}
\end{exercise}

\begin{exercise} 两个副本的 P99 分别为1秒和5秒，为什么不能直接报告全局 P99 为3秒？

\begin{solution}
全局分布是按请求数加权的混合分布，不是各分位数的平均。举例两个副本都固定分别耗1秒与5秒，只要5秒请求占比超过1\%，全局p99就是5秒；占比很小时可能为1秒。应合并同口径直方图或逐例样本，保留计数权重和桶精度。
\end{solution}
\end{exercise}

\begin{exercise} 失败概率0.2、最多四次总尝试，计算期望重试放大系数，并说明独立假设在过载时为何可疑。

\begin{solution}
总尝试次数$N$满足 $E[N]=\sum_{j=1}^4P(N\ge j)=1+0.2+0.2^2+0.2^3=1.248$，平均重试数0.248。过载故障依赖共享资源，尝试相关且重试进一步增加负载；应采用总预算、退避抖动和熔断，而非据此认定只增加24.8\%流量。
\end{solution}
\end{exercise}

\begin{exercise} 制定一次工作进程失联演练的观察项，覆盖在途输出、缓存释放、错误计量和剩余容量。

\begin{solution}
先记录故障前活动请求、页所有权与可服务副本；让指定工作进程失联，观察路由停止准入的时间、已输出前缀及唯一终态。检查页释放或迁移代次、迟到事件、重复请求、未知结果及恢复后是否双重输出。分别报告提交失败率、接纳后失败率、受影响请求恢复时长和剩余达标吞吐；不能把失联期间请求从统计分母删除。
\end{solution}
\end{exercise}

\end{exercises}
